% ════════════
% EDI_Working_Paper.tex
% Economic Dignity Infrastructure: The Agentive Reintegration Pathway
% Material Dignity Infrastructure Working Paper 5
% ════════════

\documentclass[12pt, letterpaper]{article}

% ── PACKAGES ──
\usepackage[left=0.7in, right=0.7in, top=1in, bottom=1in, headheight=14pt]{geometry}
\usepackage[expansion=false]{microtype}
\usepackage{textcomp}
\usepackage{setspace}
\usepackage{booktabs}
\usepackage{tabularx}
\usepackage[titles]{tocloft}
\usepackage{tabularray}
\usepackage{longtable}
\usepackage{array}
\usepackage{multirow}
\usepackage{hyperref}
% natbib not used — manual thebibliography
\usepackage{amsmath}
\usepackage{amssymb}
\usepackage{parskip}
\usepackage{titlesec}
\usepackage{fancyhdr}
\usepackage[table]{xcolor}
\usepackage{caption}
\usepackage{footnote}
\usepackage{enumitem}
\usepackage{tikz}
\usepackage{needspace}
\usepackage{etoolbox}
\usepackage{float}

% ── COLOR SYSTEM ───
\definecolor{auditblue}{HTML}{003366}
\definecolor{rowgray}{HTML}{F5F5F5}

% ── GLOBAL SPACING ───
\setstretch{1.4}
\setlength{\parindent}{0pt}
\setlength{\parskip}{8pt}
\hyphenpenalty=1000

% ── SECTION FORMATTING. ZERO BOLDING IN TEXT. ───
\titleformat{\section}
  {\normalfont\large\color{auditblue}}{\thesection.}{0.5em}{}
\titleformat{\subsection}
  {\normalfont\normalsize\color{auditblue}}{\thesubsection.}{0.5em}{}
\titlespacing*{\section}{0pt}{18pt}{6pt}
\titlespacing*{\subsection}{0pt}{12pt}{4pt}

% ── HEADER / FOOTER ───
\pagestyle{fancy}
\fancyhf{}
\rhead{\small\textit{{Material Dignity Infrastructure Working Paper 5}}}
\lhead{\small\textit{Working Paper, June 2026}}
\cfoot{\thepage}
\renewcommand{\headrulewidth}{0.4pt}

% ── TABLE SETTINGS ───
\renewcommand{\arraystretch}{1.3}

% ── ORPHAN / WIDOW CONTROL ───
\clubpenalty=10000
\widowpenalty=10000
\displaywidowpenalty=10000

% ── BIBLIOGRAPHY ORPHAN PREVENTION ───
\AtBeginEnvironment{thebibliography}{\interlinepenalty=10000}

% ── HYPERREF CONFIGURATION ───
\hypersetup{
  colorlinks=true,
  linkcolor=black,
  citecolor=black,
  urlcolor=black,
  pdftitle={{Economic Dignity Infrastructure}: {The Agentive Reintegration Pathway}},
  pdfauthor={{Charles J. DiBella}},
  pdfsubject={Working Paper},
  pdfkeywords={{Economic Dignity Infrastructure, agentive selfhood, identity capital, return deficit, worker cooperatives, Individual Placement and Support, CalAIM billing, Tenancy Bridge Guarantee, structural labor transitions, homelessness recovery}}
}
\setcounter{tocdepth}{2}

\begin{document}
\captionsetup[table]{labelfont=normalfont, labelsep=period, skip=6pt}

% ── FRONT MATTER ───
\begin{titlepage}
\begin{center}
\setstretch{1.4}
{\LARGE \textbf{Economic Dignity Infrastructure}}\\[6pt]
{\large The Agentive Reintegration Pathway}\\[0.5cm]
{\normalsize \textbf{Charles J. DiBella}}\\[0.01cm]
{\normalsize Principal Systems Architect}\\[0.01cm]
{\normalsize Working Paper, June 2026}\\[0.3cm]
\textbf{ABSTRACT}
\vspace{0.2cm}
\end{center}
\begin{minipage}{0.95\textwidth}
\setstretch{1.15}
\noindent

Material stabilization within a dedicated institutional environment does not, on its own, produce societal reintegration. Individuals stabilized materially and relationally still face a competitive external labor market that penalizes r\'{e}sum\'{e} gaps, institutional histories, and the absence of normative identity capital. When stabilized residents attempt traditional labor market reentry, they encounter the Return Deficit: the aggregate, compounding reduction in labor market returns produced by carrying multiple co-occurring stigmatized institutional markers simultaneously. This paper synthesizes established findings across multiple disciplines into a proposed institutional design. It develops a theoretical framework and specifies mechanisms intended for future empirical evaluation rather than presenting new causal estimates. The paper introduces the Return Deficit as its principal theoretical contribution, specifies the cooperative reintegration mechanism and the Tenancy Bridge Guarantee, and identifies the failure modes that implementation must address. Detailed financial modeling is reserved for a subsequent working paper in this series.

\end{minipage}
\end{titlepage}

\pagenumbering{arabic}

% ── DOCUMENT METADATA ────
\newpage
\section*{Document Metadata}

\noindent\textbf{Keywords}\\[1pt]
Economic Dignity Infrastructure, agentive selfhood, identity capital, return deficit, worker cooperatives, Individual Placement and Support, CalAIM billing, Tenancy Bridge Guarantee, structural labor transitions, homelessness recovery

\vspace{8pt}
\noindent\textbf{JEL Classification}
\begin{itemize}[nosep, before={\vspace{-8pt}}, leftmargin=2em]
    \item[] \textit{Primary:}
    \item I38. Government Policy; Provision and Effects of Welfare Programs
    \item J24. Human Capital; Skills; Occupational Choice; Labor Productivity
    \item J68. Labor Market Policy
    \item[] \textit{Secondary:}
    \item D63. Equity, Justice, Inequality, and Other Normative Criteria
    \item P13. Cooperative Enterprises
    \item R23. Regional Migration; Regional Labor Markets; Population; Neighborhood Characteristics
\end{itemize}

\vspace{8pt}
\noindent\textbf{Target eJournals}
\begin{itemize}[nosep, before={\vspace{-8pt}}, leftmargin=2em]
    \item Social Enterprise eJournal
\end{itemize}

\vspace{8pt}
\noindent\textbf{Author's Note}\\[1pt]
This paper is the fifth working paper in the Material Dignity Infrastructure series. MDI-D established the relational architecture necessary to produce ontological security within the physical tower. This paper establishes the economic architecture that translates that security into formal, sovereign independence, completing the full street-to-home theoretical pipeline.

\vspace{8pt}
\noindent\textbf{Suggested Citation}\\[1pt]
DiBella, C.J. (2026). Economic Dignity Infrastructure: The Agentive Reintegration Pathway. Material Dignity Infrastructure Working Paper 5, June 2026.

\newpage
% ── TABLE OF CONTENTS ────
\vspace{12pt}
\begingroup
\setstretch{1.35}
\setlength{\cftbeforesecskip}{2pt}
\setlength{\cftbeforesubsecskip}{-2pt}
\tableofcontents
\endgroup
% ── BODY ────
\setstretch{1.35}

\section{Introduction}

The preceding papers in this series defined the mechanical and relational scaffolding required to stabilize urban populations experiencing chronic physiological collapse. Material Dignity Infrastructure (MDI) established the physical environment necessary to reverse metabolic deterioration through secure acoustic sanctuaries and caloric predictability. Relational Dignity Infrastructure (RDI) subsequently mapped the social architecture, built upon Dunbar-scale cohorts and live-in Pod Stewards, required to produce the ontological security that makes physical housing inhabitable. Stabilization within a dedicated facility environment does not, on its own, equate to societal reintegration. A structural integration gap remains between facility stability and competitive labor market participation.

Durable exit from the stabilization pipeline requires a final architectural component: Economic Dignity Infrastructure (EDI). This paper synthesizes established findings across multiple disciplines into a proposed programmatic design. It develops a theoretical framework and specifies mechanisms intended for future empirical evaluation rather than presenting new causal estimates. The central theoretical anchor of EDI is Amartya Sen's capability approach \cite{sen1999development}, which frames poverty as a deprivation of basic capabilities rather than restricted to income deprivation.

To operationalize this approach, this architecture defines three core theoretical constructs, developed for this paper series. First, it defines the Return Deficit, a construct introduced in this paper to describe the aggregate, compounding reduction in labor market returns experienced by individuals carrying multiple co-occurring stigmatized institutional markers. Second, it addresses identity capital, grounded in C\^{o}t\'{e} and Levine's identity formation framework \cite{cote2002identity}, the intangible psychological currency and professional narrative required for labor market navigation. Third, it specifies the conditions for agentive selfhood, the behavioral capacity to initiate unprompted action, participate in governance, and articulate long-term economic goals.

Individuals stabilized materially and relationally still face a competitive external labor market that systematically penalizes r\'{e}sum\'{e} gaps and institutional histories. Without a formal mechanism to convert physical stability into agentive selfhood, the stabilization pipeline functions as an advanced containment system. By establishing internal cooperative economies and leveraging evidence-based Individual Placement and Support (IPS) models, EDI is designed to produce the identity capital required for transition. This paper specifies the cooperative reintegration mechanism and the Tenancy Bridge Guarantee, detailing how stabilized individuals could graduate from program support into independent housing and competitive labor market participation.

\section{The Return Deficit: A Theoretical Construct}

The architectural sequence of the MDI framework dictates that employment cannot precede biological stabilization. Neurological data confirms that chronic caloric insufficiency produces measurable changes in behavioral economics \cite{mullainathan2013scarcity}. Building upon classical starvation research \cite{keys1950biology}, subjects operating under severe resource deprivation demonstrate extreme impulsivity and an overwhelming preference for immediate over delayed rewards, rendering long-term employment planning biologically impossible. Reintegration into formal labor structures cannot occur until these neurological deficits are reversed during Phase Zero metabolic stabilization.

Once physiological stabilization is achieved, residents encounter a distinct structural barrier that existing labor market constructs describe only partially. This paper introduces the Return Deficit as a construct designed to capture this residual barrier in its entirety. The Return Deficit is defined here as the aggregate, compounding reduction in labor market returns experienced by individuals who carry multiple co-occurring stigmatized institutional markers, where each marker independently reduces employer responsiveness, and where the compound effect of carrying several simultaneously is not additive but multiplicative, producing outcomes that no single-construct explanation accounts for.

A reviewer may ask why this compound phenomenon requires a named construct rather than description as multiple intersecting disadvantages. The answer is architectural. Existing interventions optimize against individual barriers because the existing literature studies barriers individually. Each sub-field, labor scarring research, signaling theory, audit studies of criminal record penalties, has produced a well-specified single-variable finding and a corresponding single-variable intervention. When the compound interaction among these barriers becomes the unit of analysis rather than any individual component, a different class of response becomes rational: one that addresses the full marker set simultaneously rather than sequentially. Sequential intervention against a multiplicative penalty structure cannot close the deficit; it reduces one factor while the others remain active. Naming the compound phenomenon as the target of analysis is the necessary precondition for specifying a response calibrated to its actual structure.

Labor economists have documented several related phenomena independently. Labor market scarring refers to the permanent wage and employment penalties that attach to workers following extended unemployment spells, with effects that persist years beyond the initial spell's end \cite{pager2003mark}. Signaling theory explains how employers use observable credentials as proxies for unobservable productivity, penalizing workers whose signals have degraded. Statistical discrimination describes how employers apply group-level probability assessments to individual candidates in the absence of direct information, producing systematic exclusion that is rational from the firm's perspective and destructive from the worker's. Criminal record penalties, documented extensively through audit studies using matched r\'{e}sum\'{e}s, produce callback rate reductions of thirty to fifty percent for otherwise identical applicants. R\'{e}sum\'{e} gap penalties operate as a separate signal of disengagement, producing additional independent reductions in employer responsiveness even when the gap's cause is entirely outside the worker's control.

These constructs share a common assumption: the penalty under study is the primary or dominant barrier facing the worker being analyzed. An individual carrying a single r\'{e}sum\'{e} gap, a single criminal record, or a single period of unemployment faces a measurable deficit that existing interventions, vocational training, certificate programs, and transitional employment, are designed to address. The individual exiting long-term unsheltered homelessness does not present a single stigmatized marker. That individual presents, simultaneously, an extended employment gap measured in years rather than months, the absence of a fixed address during the gap period, the absence of professional references capable of verifying any work history during the gap period, a probable criminal justice intersection, and in many cases, a visible institutional history traceable through public records and background screening services. Each marker independently produces the penalties documented in the existing literature. Carried together, they produce a compounding return structure that pushes effective labor market returns toward zero regardless of the effort invested.

The Return Deficit names this compounding structure as a single institutional phenomenon rather than a list of co-occurring penalties. This naming matters for institutional design because an intervention calibrated to address one component of the deficit leaves the remaining components intact and operative. A transitional employment program that fills the r\'{e}sum\'{e} gap leaves the criminal record penalty untouched. A certificate program that adds a credential leaves the address history gap untouched. A supported employment placement that creates a professional reference leaves the statistical discrimination operating against the individual's demographic profile untouched. The Return Deficit framework predicts that incremental interventions will consistently underperform their designers' expectations, producing short-term placements followed by recidivism into institutional dependence, not because of individual failure but because the structural friction has been reduced rather than eliminated.

The Return Deficit is measurable. Its primary indicators are employer callback rates for matched r\'{e}sum\'{e} profiles across stigmatized and non-stigmatized cohorts, employment spell duration in the twelve and twenty-four months following placement, wage trajectory relative to non-stigmatized entrants at equivalent skill levels, and the rate at which placed workers return to institutional support within eighteen months. Secondary indicators include the frequency with which initially placed workers report employer disengagement, schedule reduction, or involuntary termination at the first sign of the institutional history becoming visible during employment. These are not new data collection requirements. They are existing metrics measured in isolation by existing research programs, which the Return Deficit framework synthesizes into a single longitudinal index.

Traditional vocational rehabilitation and Employment Social Enterprises address components of this deficit sequentially and incompletely. Housing First methodologies \cite{tsemberis2004housing} resolve biological stabilization and residential address history effectively but leave the employment-side friction structure intact. Employment Social Enterprises provide wage-paying transitional labor and fill the r\'{e}sum\'{e} gap, but they return workers to the competitive external market before the compounding stigma structure has been addressed at the institutional level. This premature exposure reactivates the statistical discrimination and signaling penalties that temporary employment alone cannot neutralize. EDI is designed as an institutional response to the deficit in its compound form rather than to any single component, and the remainder of this paper specifies the mechanisms through which that response operates.

\section{Identity Capital and Agentive Selfhood Accumulation}

The transition from passive service recipient to active economic participant requires more than physical shelter; it requires the accumulation of identity capital and the restoration of agentive selfhood, as defined in Relational Dignity Infrastructure (RDI). Sociological models define identity capital as the intangible psychological currency, comprising self-efficacy, a coherent professional narrative, and agentic capacity, required for successful labor market navigation \cite{cote2002identity}. This concept runs parallel to "recovery capital" \cite{white2008recovery}, encompassing the internal and external resources necessary to sustain stabilization. Operationally, identity capital accumulation is measured using validated psychometric instruments, such as the Generalized Self-Efficacy Scale \cite{schwarzer1995generalized}, alongside defined composite indices. Agentive selfhood is measured through observable behavioral indicators, including unprompted task initiation, participation in cooperative governance, and the articulation of long-term economic goals. Marginalized individuals enter the stabilization pipeline lacking this currency, possessing stigmatized identities formed through chronic exclusion.

\subsection{Physiological Grounding for Relational Continuity}

The accumulation of identity capital and the restoration of agentive selfhood require a stable physiological and relational foundation, as established in Relational Dignity Infrastructure. Sustained exposure to unpredictable street conditions degrades the executive capacities required for long-term planning \cite{arnsten2009stress}, frequently produces substance use patterns consistent with the self-medication hypothesis \cite{khantzian1985self}, and produces neurological changes that entrench short-term priority structures at the expense of future-oriented decision-making \cite{volkow2016neurobiologic}. Scattered-site housing frameworks demonstrate substantial residential stability outcomes across broad unhoused demographics \cite{tsemberis2004housing}, and nothing here contradicts that evidence base. A distinct consideration applies to the specific subgroup of individuals transitioning directly from long-term encampment residence with deep reliance on street-based kinship networks. This paper hypothesizes that for this specific subgroup, spatial isolation from an established peer-support network requires active psychological readaptation, as the sudden disruption of street-based kinship ties creates acute relational instability. Where this dynamic is present, congregate or co-located housing models that intentionally preserve peer-support structures during the transition period may offer a targeted advantage, without displacing scattered-site placement as the default, evidence-supported model for the broader population. The physiological stabilization established in MDI is the precondition for identity capital accumulation, and the relational architecture established in RDI is the medium through which that accumulation becomes possible.

The MDI tower provides the protected environment necessary for residents to accumulate identity capital before attempting external transitions. Accumulating this capital requires structural scaffolding rather than motivational incentives. Individuals must engage in meaningful, progressively responsible actions within a protected facility ecosystem to verify their shifting self-concept through actual economic performance. Because the competitive external market penalizes early missteps during this formation period, the economic participation must occur internally, within a structural environment that captures the learning cost rather than externalizing it onto the worker.

\section{The Cooperative Reintegration Mechanism}

Operationalizing internal micro-economies requires a structural departure from traditional vocational rehabilitation models. Worker cooperatives present a potential architectural solution. Within the broader literature on hybrid organizations \cite{battilana2014advancing}, EDI models function closer to European social cooperatives \cite{defourny2010conceptions, galera2009social} rather than traditional American Employment Social Enterprises (ESEs). Unlike standard ESEs which primarily prepare individuals for eventual insertion into competitive external markets, cooperative hybrid models transition marginalized populations directly into participatory ownership structures \cite{perotin2014what}. This design ensures that every dollar spent maintaining the physical infrastructure simultaneously funds the economic reintegration of its residents. The physical tower becomes an economic engine, capturing capital that would otherwise exit the stabilization pipeline.

A facility operating at one hundred enrolled residents generates an estimated twenty-five to thirty-five compensated cooperative positions across the four primary functions. Culinary operations require three shifts of three to four workers each, producing approximately ten to twelve positions. Facility maintenance requires two to three workers per shift. Security requires two workers across three daily shifts for continuous coverage. Administrative support generates two to four positions depending on enrollment and billing complexity. This ceiling is a design constraint, not a program failure. Enrollment levels that exceed available cooperative positions require supplementary participation pathways: structured volunteer roles in resident governance, satellite cooperative operations in adjacent facilities, and graduated external supported employment placements that introduce market exposure progressively rather than abruptly.

Cooperative models require rigorous governance to prevent failure. Transaction cost economics demonstrates that cooperative ownership structures experience significant hidden costs resolving disputes among members \cite{hansmann1996ownership}. Literature on worker cooperatives identifies acute vulnerabilities, including the ``free-rider'' problem \cite{ostrom1990governing} and generative conflict fatigue, where members benefit from collective success without contributing proportional labor. The MDI cooperative model applies Ostrom's core design principles directly: membership boundaries are defined by residency status, decision-making rules are established through bylaws adopted collectively before operations begin, contributions are monitored through Pod Steward-maintained labor logs, and sanctions escalate from peer mediation to facility management review. The profit distribution formula links total hours logged to distribution share, with a governance participation multiplier that adds a percentage premium for attendance at cooperative meetings, preventing the concentration of earnings among a small group of high-tenure members. Elected cooperative roles carry fixed term limits to prevent capture by founding members and require mandatory rotation at defined intervals.

To maximize retention and ensure functional competence, these internal cooperatives integrate the evidence-based Individual Placement and Support (IPS) model \cite{drake2012individual}. Operating on a "place, then train" principle, IPS eliminates prolonged job readiness prerequisites that frequently discourage marginalized participants. A recent clinical trial examining IPS-supported employment among justice-involved homeless veterans demonstrated superiority over standard vocational rehabilitation when integrated with clinical support \cite{lepage2021individualized}. By embedding IPS protocols directly within the cooperative ownership structure, the MDI framework is designed to ensure that residents experience immediate economic enfranchisement while receiving the continuous clinical support necessary to sustain their performance.

\section{The Tenancy Bridge Guarantee}\label{sec:tbg}

The final structural requirement of the MDI framework addresses the exit velocity necessary to prevent permanent institutionalization. Internal cooperative micro-economies resolve the immediate identity capital deficit, but they must function as transitional accelerators rather than permanent destinations. To convert facility stability into sovereign independence, the system deploys the Tenancy Bridge Guarantee.

The Tenancy Bridge Guarantee functions as a self-executing contractual mechanism bridging Phase One stabilization (the MDI tower) and Phase Two sovereign tenancy. Once a resident demonstrates verified identity capital accumulation, measured through sustained participation within the internal cooperative micro-economy over a designated multi-month window, the mechanism triggers autonomous capital routing. The revenue stream funding this mechanism is not hypothetical. California's CalAIM program, administered by the Department of Health Care Services, designates Housing Transition Navigation Services, Housing Tenancy and Sustaining Services, and Enhanced Care Management as separately reimbursable Community Supports under managed care plan contracting. Reimbursement rates for these categories are negotiated between providers and individual Managed Care Plans (MCPs), rather than set at a fixed statewide rate, making the revenue stream real but locally variable. Analogous billing categories exist in Oregon's Coordinated Care Organization framework and in Washington State's Medicaid waiver program, establishing that the model is not dependent on a single state's policy architecture. While federal Medicaid funding frequently prohibits direct room and board payments, the clinical revenue produced by housing navigation and care management services underwrites the administrative overhead required to secure and guarantee external market-rate master leases.

The financial architecture supporting the Tenancy Bridge Guarantee operates across three revenue sources and three cost centers, specified here at the categorical level. Detailed financial modeling, including unit-level projections, capitation rate assumptions, and operating gap analysis, is reserved for a subsequent working paper in this series.

Revenue sources supporting the guarantee are: first, the cooperative enterprise operations, which internalize labor costs that would otherwise flow to external vendors, reducing net facility expenditure while generating employment income for residents; second, Medicaid care management navigation reimbursement, available through managed care plan contracting for housing transition and enhanced care management services, which the managing entity earns by maintaining enrolled residents through the stabilization-to-tenancy transition; and third, administrative support fees associated with master lease negotiation and ongoing tenancy support for graduated residents. The primary cost centers are facility operations, clinical and stewardship staffing, and the transition support services required to secure and guarantee external market-rate leases on behalf of graduating residents.

The financial logic connecting these elements is as follows. The cooperative enterprises reduce the facility's labor cost outlay, generating a margin that does not exist in conventional shelter operations. The Medicaid navigation revenue converts the clinical work of stabilizing and transitioning residents into a recurring revenue stream rather than a grant-dependent expenditure. Together, these sources are designed to reduce the operating gap that would otherwise require full philanthropic or public subsidy, without the model claiming that this reduction fully closes that gap under all operating conditions. The scale and timing of the gap reduction are empirical questions dependent on local Medicaid waiver structures, cooperative revenue generation, and enrollment stability, and they are not resolved by this paper.

Residents who complete the Tenancy Bridge Guarantee exit with durable income histories produced by the cooperative, professional references from cooperative management and Pod Stewards, and ongoing relational support from the network established during stabilization.

\section{Theory of Change}

The theoretical claims underpinning Economic Dignity Infrastructure require rigorous empirical testing before widespread adoption. Standard programmatic metrics, including immediate job placement rates, do not capture the structural integration EDI proposes. Evaluation must follow an explicit causal pathway connecting specific programmatic mechanisms to measurable longitudinal outcomes.

The EDI theory of change proceeds as follows. Street instability produces physiological and neurological conditions that preclude sustained labor participation. Phase Zero stabilization, delivered through MDI, reverses the acute physiological deficit. Relational stability, established through RDI's cohort and Pod Steward architecture, creates the social foundation required for identity capital accumulation. Protected economic participation, within the internal cooperative economy, produces the first empirical evidence of agentive selfhood and begins accumulating the income history, professional reference structure, and psychological capital that constitute a measurable reduction in the Return Deficit. Sufficient capital accumulation triggers the Tenancy Bridge Guarantee, routing program revenue into external tenancy support and enabling competitive labor market participation from a position of residential stability rather than chronic precarity.

Each transition in this sequence is a hypothesis, not a demonstrated outcome. The construct validity of the Return Deficit as a unified index, the degree to which internal cooperative participation produces transferable identity capital, the reimbursement stability of Medicaid navigation revenue, and the recidivism rate among Tenancy Bridge Guarantee graduates are all empirical questions that future evaluation must address. This paper establishes the governance logic and specifies the mechanisms; it does not validate them.

Primary evaluation metrics include sustained external housing retention at twenty-four and thirty-six months following placement, income trajectory in the twelve months after cooperative-to-external transition, and emergency service utilization rates across the stabilization and post-graduation periods. Secondary psychometric evaluation must track identity capital accumulation using validated instruments across the full residency period to test the hypothesis that internal cooperative participation restores agentive selfhood before competitive external market exposure.

A pilot implementation cannot wait twenty-four months for primary outcomes before assessing whether the mechanism is functioning. The following leading indicators are observable within the first twelve months of operation and are derivable from payroll records, Pod Steward logs, cooperative meeting attendance sheets, and the psychometric instruments already embedded in the framework. Monthly indicators include: cooperative labor participation rate, defined as the proportion of enrolled residents holding an active compensated position; governance meeting attendance rate; and unprompted task initiation count, logged by Pod Stewards against a defined behavioral protocol. Quarterly indicators include Generalized Self-Efficacy Scale scores and accumulated internal wage hours per resident. At six months, a binary professional reference indicator records whether each resident has at least one verifiable reference available for external job applications. At twelve months, the external application rate and callback rate for residents who have begun external job-seeking activity provide the first employment-side signal that the Return Deficit is measurable and responding to the intervention. These indicators do not replace the primary outcomes. They are soil moisture readings taken before the harvest to determine whether conditions are supporting growth.

\section{Failure Modes and Structural Limitations}

System-level design proposals carry an obligation to specify the conditions under which the proposed mechanisms fail. The EDI framework contains several failure modes that are not peripheral risks but predictable structural vulnerabilities that implementation must address directly.

Cooperative governance failure represents the primary internal threat. The worker cooperative literature documents persistent difficulties in sustaining democratic governance under resource pressure, including the concentration of decision-making authority in founding members, the erosion of bylaws during financial stress, and the displacement of member interests by organizational survival demands. The EDI model's reliance on resident governance participation as a mechanism for identity capital accumulation presupposes that governance structures remain functional under these conditions. Where they do not, the cooperative becomes a managed employment program with ownership language attached, producing neither the organizational loyalty nor the agentive selfhood accumulation the model requires.

The free-rider problem is structurally embedded in any cooperative arrangement linking individual benefit to collective performance. Residents who receive housing stabilization and clinical support from the MDI structure have an incentive to remain enrolled regardless of their labor contribution to the cooperative. Formal dispute resolution bylaws and contribution-tracked profit distribution reduce this risk but do not eliminate it. A sufficiently high proportion of low-contributing members degrades cooperative revenue, reduces the margin available to underwrite the Tenancy Bridge Guarantee, and produces organizational conflict that consumes the governance capacity the model depends upon.

Medicaid reimbursement policy represents an external structural vulnerability that the model cannot internally address. The revenue architecture described in Section~\ref{sec:tbg} depends on managed care plan contracting for housing transition navigation services, and on the continued availability of enhanced care management billing categories at rates sufficient to produce a financial margin. These conditions reflect the current state of Medicaid waiver policy in states that have adopted relevant managed care provisions, and they are not uniform across jurisdictions, nor are they stable over time. A reimbursement rate reduction, a billing category reclassification, or a managed care plan contracting decision can remove a primary revenue source without any corresponding reduction in facility costs.

Insufficient internal labor demand is a scaling constraint that becomes critical as enrollment increases. The cooperative enterprises described, maintenance, culinary operations, security, and administrative support, produce a finite quantity of compensated labor positions within a single facility. An enrollment level that exceeds available cooperative labor positions creates a cohort of residents who are stabilized and enrolled but unable to participate in the primary identity capital accumulation mechanism. Supplementary mechanisms, including structured volunteer roles, external supported employment placements, or satellite cooperative operations, would be required to maintain the participation pathway at scale, and each introduces the external market exposure that the model is designed to defer.

Participant heterogeneity poses a challenge to any fixed-sequence programmatic framework. The MDI, RDI, and EDI progression assumes that physiological stabilization, relational attachment, and economic participation follow a predictable developmental sequence applicable across a diverse enrolled population. For some residents, the sequence will compress; for others, the timeline will extend significantly beyond the parameters on which the financial model is built. Residents with severe co-occurring disorders, complex trauma histories, or acute medical needs may require clinical resources that exceed what the cooperative revenue structure can support, and may occupy stabilization capacity for periods inconsistent with the model's transition assumptions.

Permanent institutionalization risk, the possibility that the MDI structure becomes a terminal destination rather than a transitional mechanism, functions as a structural design problem rather than a clinical failure. A resident who remains enrolled beyond the model's intended transition window consumes stabilization resources, occupies a cooperative labor position, and fails to produce the Tenancy Bridge revenue that the financial architecture requires to remain solvent. Explicit residency timelines, transition accountability structures, and discharge planning protocols must be specified at the facility level to prevent the stabilization pipeline from accumulating a permanent census rather than producing a continuous throughput of graduating residents.

These failure modes do not invalidate the framework. They specify the conditions under which the framework's performance will degrade and identify the design choices that implementation must address before the model can be evaluated fairly against the outcomes it proposes to produce.

\section{Conclusion}

Economic Dignity Infrastructure completes the three-stage structural framework. Material housing addresses the biological preconditions for stable function. Relational infrastructure produces the social architecture within which identity capital can accumulate. The internal cooperative economy provides the protected structural mechanism through which that capital becomes measurable and transferable. The framework hypothesizes that chronic unsheltered homelessness is not resolved by returning stabilized individuals to the competitive external markets that produced their exclusion before the compound Return Deficit has been addressed at the structural level.

Working Paper 5 functions as the conceptual bridge between the architectural papers that precede it and the empirical validation work that must follow. Future papers in this series will address financial modeling, evaluation design, and implementation analysis at the site level. The present paper's contribution is the Return Deficit construct and the structural architecture specified to address it.

\newpage
\addcontentsline{toc}{section}{References}
\bibliographystyle{apalike}
\bibliography{../01_drafts/bibliography_final}

\end{document}
